% Cover letter: AI Engineer, Parallel (Paris, on-site), 2026-09-23.
% Built with its CV profile of the same name, so the header carries the
% parallel-ai tagline and "Open to relocating to Paris":
%   make letter-parallel-ai   ->   Cover_Letter_<First>_<Last>_Parallel_AI.pdf
%
% Every claim is on the parallel-ai CV or was confirmed by the owner in the
% session that wrote it: the self-labeling tool (-75% hand-labeling), the
% coding agent (co-started, production, self-hosted GLM-5.3 with vision
% capabilities from Hugging Face, 6 H200s, ~10x throughput, NVFP4 + FP8 KV
% cache, 1M-token context, privacy-first), the AI gateway (Go WebAssembly
% plugins, 45 researchers), more than four industrial vision projects owned
% end to end, ~14 fps on the Hailo-fitted PLC, the customer field-feedback
% loop, RemaNet (25 partners; React, TypeScript, Node.js), MLflow/ClearML,
% CI/CD and GPU cluster with Grafana/Prometheus, the AI-first daily
% practice, French B1 and English C1.
%
% Deliberately not claimed: healthcare or medical-coding experience,
% regulated-data compliance, LangChain. The letter names healthcare as new
% ground rather than leaving the reader to find the gap.
%
% The opening picks up the posting's own first line, "up to 25% of
% healthcare spending is lost in manual, repetitive administrative work"
% (owner's suggestion, 2026-09-23).
%
% Salutation: "Dear Parallel team". The posting names the recruiter
% (Wallerand) and the CTO (Chris); "Dear Wallerand, dear Chris" is the more
% personal alternative. Kept here, not rendered:
%   Dear Wallerand, dear Chris,
%
% Nothing in this file is ever deleted: a replaced sentence goes into a
% comment next to its replacement, as in the CV sources.

\renewcommand{\cvfootertitle}{\IfLanguageName{english}{Cover Letter}{Anschreiben}}

% ---------------------------------------------------------------------------
% Version 1 (2026-09-23), kept, not rendered. The owner's verdict: "not just a
% repetition of what is on the CV". It walked the CV's facts paragraph by
% paragraph. Version 2 below keeps the facts but tells three lessons instead,
% each answering a line of the posting.
%
% \renewcommand{\cvBody}{%
% 	\cvletterhead
% 		{Stuttgart, 23 September 2026}
% 		{Parallel\\AI Engineer hiring team\\Paris, France}
% 		{Application: AI Engineer (Paris, on-site)}
%
% 	Dear Parallel team,
%
% 	Up to 25\% of healthcare spending lost in manual, repetitive administrative work: that describes the kind of problem I have spent the last two and a half years solving at Fraunhofer IPA, in another domain. My self-labeling tool built on vision-language models cut hand-labeling time by 75\%, and the coding agent I started with a colleague now takes repetitive work off the institute's researchers as a production service. I would like to bring that experience to hospitals as an AI Engineer at Parallel.
%
% 	The agent is an LLM system that runs, not a proof of concept: a self-hosted, vision-capable GLM-5.3 with weights from Hugging Face on six H200 GPUs, made fast enough for daily use with custom speculative decoding and scheduling ($\sim$10$\times$ throughput) and an NVFP4 model with an FP8 KV cache for a 1M-token context. Around it, I extended the institute's AI gateway with WebAssembly plugins in Go for token metering and rate limiting across 45 researchers. We run all of it privacy-first on our own hardware, which is the same discipline health data demands of your pipelines.
%
% 	I own what I build end to end. In more than four industrial vision projects I built and annotated the datasets, trained the models in PyTorch, evaluated and deployed them, down to an embedded PLC at $\sim$14\,fps, and iterated with the industry customers on their field feedback. For RemaNet, a 25-partner Horizon Europe project, that included the React frontend (TypeScript, Node.js) through which people use the results. Experiments are tracked in MLflow and ClearML; deployments go through the GitLab CI/CD and GPU cluster I run, monitored with Grafana and Prometheus.
%
% 	I already work the way your posting describes: coding agents, AI-assisted review and AI-driven testing and debugging are my default tools. Healthcare and medical coding would be new ground for me; automating repetitive expert work with AI systems that people rely on every day is not. I am glad to move to Paris for an on-site role, and I speak French at B1 and English at C1.
%
% 	I would welcome the chance to discuss how I could help build Parallel's AI systems.
%
% 	\cvletterclosing{Best regards}
% }
% ---------------------------------------------------------------------------
%
% Version 2 (2026-09-23): three lessons, each tied to a line of the posting.
%   hook      -> the posting's own first line, and what the owner has already
%                removed of that kind of work (75% of hand-labeling; the agent)
%   lesson 1  -> "real systems that run, not just POCs": what took the agent
%                from side project to production was speed and privacy work
%   lesson 2  -> "No integrations, no disruption": the retrofitted Berghof PLC
%                and the customer field-feedback loop
%   lesson 3  -> "sensitive health data" and "engineering rigor ... automated
%                best practices ... tools like Cursor": privacy-first by policy,
%                and this letter itself -- typeset by the owner's open-source
%                pipeline with tests and strict gates, written with a coding
%                agent (true of this letter; the footer links the pipeline)
%   gap       -> medical coding named as the one thing not brought yet, and
%                how it would be learned
% Sentences that state what the owner learned or thinks ("taught me", "for
% me") are drafted from the owner's facts and wording in the session. The
% owner should read them as a proposal and put them in their own voice where
% they differ. HDS (France's certification for hosting health data) is named
% as the bar to clear, not as experience.
\renewcommand{\cvBody}{%
	\cvletterhead
		{Stuttgart, 23 September 2026}
		{Parallel\\AI Engineer hiring team\\Paris, France}
		{Application: AI Engineer (Paris, on-site)}

	Dear Parallel team,

	% First wording of the second sentence, kept ("work" three times in two
	% sentences): I know that problem from the builder's side: much of my work at Fraunhofer IPA removes exactly that kind of work.
	Your posting opens with up to 25\% of healthcare spending lost to manual, repetitive work. I know that problem from the builder's side: much of what I build at Fraunhofer IPA removes exactly that kind of work. A self-labeling tool I built on vision-language models cut hand-labeling time by 75\%, and a coding agent that a colleague and I started on our own initiative is now a production service the institute's researchers use. Three things I learned doing it are why I am writing to you.

	\textbf{People only hand work to an agent that is fast and trustworthy.} What took our agent from side project to production was the unglamorous part: custom speculative decoding and scheduling for $\sim$10$\times$ throughput, NVFP4 with an FP8 KV cache for a 1M-token context, a self-hosted GLM-5.3 from Hugging Face, served with vLLM, so our code stays in-house, and token metering and rate limiting in the institute's AI gateway.

	\textbf{The best integration is none.} Your ``no integrations, no disruption'' has an industrial twin in my work: I retrofitted a Berghof PLC with a Hailo AI chip, ran PyTorch-trained, INT8-quantized object detection on it at $\sim$14\,fps, and then kept iterating with our industry customers on their field feedback.

	\textbf{Rigor should be automated.} We self-host because research code should not leave the institute; with patient data and HDS-grade hosting, that constraint becomes the product, and I would like to learn to meet it with you. Even this letter was typeset by my open-source document pipeline, with tests and strict build gates, and written the way I write code: with a coding agent at my side.

	What I do not bring yet is medical coding. I would learn it the way I learned remanufacturing for RemaNet and surface treatment for KI-PrOT: close to the people who do the work, iterating on their feedback until the system earns their trust. I would be glad to move to Paris for this role, and I speak French at B1 and English at C1.

	I look forward to showing you in the case study how I would approach your first problem.

	\cvletterclosing{Best regards}
}
